\section{Relative tolerance, tensor blocks, and escape cuts}
\label{sec:relative-blocks}
Many independent blocks of fixed size give a fixed relative-gap
obstruction for increasing concave costs. The node oracle is still
global: its square multipliers can couple variables from every block.
The blocks are independent in the optimization problem, not in the
oracle; this independence also gives a short alternative certificate
outside the single-cover model.

Fix integers $r\ge1$, $t\ge2r-1$, and $G\ge1$. There are $G$ blocks,
each with $\ell=3t$ original coordinates and its own balance:
\begin{equation}\label{eq:relative-problem}
 \begin{split}
 n&=3tG,\qquad K=t+\tfrac12,\qquad
 \sum_{i=1}^{3t}x_{b,i}=K\quad(b=1,\ldots,G),\qquad 0\le x_{b,i}\le1,\\
 C(x)&=\sum_{b=1}^G\sum_{i=1}^{3t}
       \bigl[x_{b,i}(2-x_{b,i})+d_{b,i}x_{b,i}\bigr].
 \end{split}
\end{equation}
Assume $d_{b,1}<\cdots<d_{b,3t}$ are positive and
$\sum_i d_{b,i}\le\eta$ for every block, with $\eta>0$.
Every coordinate cost is strictly increasing and strictly concave, by
its derivatives $2-2x+d_{b,i}>0$ and $-2$.
Corollary~\ref{cor:unique-cardinality} applied blockwise gives the exact
optimum
\begin{equation}\label{eq:relative-optimum}
 C^*=G(K+\tfrac14)+
       \sum_{b=1}^G\left(\sum_{i\le t}d_{b,i}+\tfrac12d_{b,t+1}\right)
 \ge G(K+\tfrac14).
\end{equation}

On an arbitrary coordinate product domain, grant the order-$r$ oracle
\eqref{eq:product-domain-oracle}--\eqref{eq:product-domain-equalities}
with all $G$ balance equalities and every polynomial multiplier through
\emph{total} degree $2r$. Squares may involve all $n$ variables. This
includes the box preordering and, for boxes at $r=1$, full global
SDP--RLT with every balance product. Let $\operatorname{LB}_r(S)$ be
the infimum of $L[C]$. A relative certificate at tolerance
$0<\theta<1$ is a finite cover of the feasible set by regions that are
infeasible or satisfy
\begin{equation}\label{eq:relative-target}
 \operatorname{LB}_r(S)\ge T_*=(1-\theta)C^*.
\end{equation}
Pruning at $(1-\theta)U$ for any feasible incumbent $U\ge C^*$ implies
this condition. The dimension in all bounds below is the original
dimension $n=3tG$.

\begin{lemma}[Tensor positivity with a total-degree truncation]
\label{lem:tensor-preordering}
For each block $b$, suppose $L_b$ is a normalized functional on block
polynomials through degree $2r$, satisfying the block balance products,
local equalities, and every block preordering inequality. Define $L$ on
global monomials of degree at most $2r$ by
\begin{equation}\label{eq:tensor-moment}
 L\left[\prod_b x_b^{\alpha_b}\right]=
             \prod_b L_b[x_b^{\alpha_b}],
\end{equation}
and extend linearly. Then $L$ satisfies the corresponding global oracle,
including every globally coupled square multiplier.
The same assertion holds when a block contains coordinatewise graph
auxiliaries and degree is measured in the lifted variables.
\end{lemma}
\begin{proof}
Each block degree is at most the global degree, so the definition is
well-defined and normalized. To test a balance equality from block $b$
times a global monomial, factor the latter across blocks. Its expectation
is the block-$b$ balance product, of degree at most $2r$, times moments
of the other blocks, so it vanishes. Linearity treats every multiplier.
The argument for local equalities is identical and also covers products
involving several equalities within the total degree bound.

For a preordering expression let $g=\prod_b g_b$ group all local
inequality factors according to their block, and let $d=\deg P$ with
$\sum_b\deg g_b+2d\le2r$. For every $b$, form the matrix
\[
 M_b(\alpha,\beta)=L_b[g_b x_b^\alpha x_b^\beta],
 \qquad |\alpha|,|\beta|\le d.
\]
It is defined and positive semidefinite: the degree bound
$\deg g_b+2d\le2r$ allows the block preordering inequality for every
polynomial in this index span. Their tensor product is positive
semidefinite, for example by tensoring Gram representations.
Keep only tuples of indices whose total degree across blocks is at most
$d$. The corresponding principal submatrix has entries exactly
$L[gx^\alpha x^\beta]$ by \eqref{eq:tensor-moment}; each such entry has
available global degree at most $2r$. Its quadratic form on the
coefficient vector of $P$ is $L[gP^2]\ge0$.
The other entries of the tensor matrix need not be globally available
moments; only the principal submatrix above is interpreted. The proof
never uses that the coordinates within a block are unlifted.
\end{proof}

\begin{theorem}[A fixed relative-gap cover bound]
\label{thm:relative-cover}
For \eqref{eq:relative-problem} and the global product-domain oracle
above, set
\begin{equation}\label{eq:relative-parameters}
 q_0=t-2r+2\ge1,\qquad
 \tau=\tfrac14-\theta(K+\tfrac14)-\eta.
\end{equation}
If $\tau>0$, every certificate at target
\eqref{eq:relative-target} has at least
\begin{equation}\label{eq:relative-cover-bound}
 (3/2)^{q_0G\tau}=(3/2)^{q_0\tau n/(3t)}
\end{equation}
regions. For fixed $r,t,\theta,\eta$ satisfying these conditions, this
is exponential in $n$. If $\tau\le0$, only the trivial bound one is
asserted. The same conclusions hold for coordinatewise graph lifts
satisfying the local oracle and full-graph objective-agreement assumptions
of Theorem~\ref{thm:local-graph-lift}, retaining all original balances;
there is no auxiliary-degree loss in $r$.
\end{theorem}
\begin{proof}
Choose independently in each block a uniform ordered partition
$(H_b,M_b,Z_b)$ into three sets of size $t$, and set
$w_b=\ind_{H_b}+p\ind_{M_b}$ with $p=1/(2t)$.
Every combined witness is feasible. Fix a product domain containing one
and define $A_b,D_b,R_b$ by endpoint exclusion within block $b$.
If $|R_b|<q_0$, the construction of
Theorem~\ref{thm:product-domain-cover} gives a feasible block functional
whose penalty is at most
\[
 |R_b|p(1-p)\le\frac{|R_b|}{2t}
                    \le\frac{|R_b|}{2q_0},
\]
since $1\le q_0\le t$. The construction uses only the endpoint
restriction count, not an absolute pruning target.
If $|R_b|\ge q_0$, instead use exact evaluation at the actual block
witness. Its penalty is
$tp(1-p)=1/2-1/(4t)\le1/2\le |R_b|/(2q_0)$.
Thus in either case the block functional is feasible through degree
$2r$, its first moments lie in $[0,1]$, and its penalty has the same
upper bound. Lemma~\ref{lem:tensor-preordering} gives a feasible global
functional of objective at most
\begin{equation}\label{eq:relative-witness-bound}
 GK+\frac1{2q_0}\sum_b|R_b|+G\eta.
\end{equation}
A certified region containing this witness must therefore satisfy
\[
 GK+\frac1{2q_0}\sum_b|R_b|+G\eta
 \ge\operatorname{LB}_r(S)\ge(1-\theta)C^*
 \ge(1-\theta)G(K+\tfrac14).
\]
It follows that $\sum_b|R_b|\ge2q_0G\tau$.

For a fixed domain, witness membership factors across blocks. In block
$b$, endpoint avoidance bounds its witness fraction by
\[
 \min\{(2/3)^{|A_b|},(2/3)^{|D_b|}\}
 \le(2/3)^{|R_b|/2}.
\]
Multiplying over independent blocks bounds the full witness fraction in
a certified region by $(2/3)^{q_0G\tau}$. The union bound over a cover
proves \eqref{eq:relative-cover-bound}. The restriction counts belong
to the region, not to the selected witness, so this is a uniform bound
for every certified region with at least one witness.

For lifted regions work with the exact original-coordinate preimages.
In a lightly restricted block use the affine endpoint substitution
\eqref{eq:graph-interpolation-map}; in a heavily restricted block use
evaluation at its true graph witness. These block functionals satisfy all
lifted localizers and equality products by the proof of
Theorem~\ref{thm:local-graph-lift}, with the original order $r$.
Lemma~\ref{lem:tensor-preordering} proves every global lifted localizer,
including squares coupling auxiliaries from all blocks. It also proves
all balance and local equality products.

A lifted polynomial objective agreeing with $C$ on the full product
graph need not be written as a sum of block polynomials.
After the block substitutions, it agrees with the substituted $C$ at
every Boolean assignment of all remaining formal variables; the heavily
restricted blocks are actual fixed graph tuples. Affine substitution
preserves its total degree at most $2r$, and uniqueness of Boolean
reduction shows the two polynomials have the same value under the tensor
functional. That functional annihilates Boolean identities in every
remaining coordinate because each factor does. Hence
\eqref{eq:relative-witness-bound} still holds, and the preimage witness
count is identical. This proves the lifted assertion without assuming
blockwise separability of the written lifted objective.
\end{proof}

For the explicit parameters
\begin{equation}\label{eq:relative-example}
 r=1,\quad t=2,\quad\theta=1/32,\quad\eta=1/32,
 \qquad K=5/2,\quad q_0=2,\quad\tau=17/128,
\end{equation}
the lower bound is $(3/2)^{17n/384}$, even for full global SDP--RLT.
For any other fixed order choose fixed $t\ge2r-1$, $0<\eta<1/4$,
and $0<\theta<(1/4-\eta)/(K+1/4)$.
This fixed-block construction gives no common positive relative
tolerance as $r$ grows without bound.

\begin{proposition}[Uniqueness and symmetry modulo the balances]
\label{prop:relative-asymmetry}
Problem~\eqref{eq:relative-problem} has a unique optimizer: in every
block its first $t$ coordinates equal one, coordinate $t+1$ equals
$1/2$, and the rest equal zero. For the explicit choice
\begin{equation}\label{eq:squared-perturbation}
 d_{b,i}=\frac{\eta\bigl((b-1)3t+i\bigr)^2}{3t\,n^2},
 \qquad b\in[G],\ i\in[3t],
\end{equation}
no nonidentity permutation of the original coordinates preserves both
the feasible set and the objective restricted to that set.
\end{proposition}
\begin{proof}
Each independent block has the unique optimizer from
Corollary~\ref{cor:unique-cardinality}; their concatenation is therefore
the unique global optimizer. The coefficients in
\eqref{eq:squared-perturbation} are positive, globally distinct, and
increasing in each block. Each is at most $\eta/(3t)$, so its block
sum is at most $\eta$ as required.

The point with every coordinate $K/(3t)$ is strictly between zero and
one. Thus the affine hull of the feasible set is exactly the displayed
balance space. Its normal row space consists of vectors constant on each
block. A feasible-set-preserving coordinate permutation preserves this
row space. The image of a block indicator is a zero-one vector in it
with precisely $3t$ ones, so it is the indicator of one entire block.
Hence such a permutation maps whole blocks to whole blocks.

The unperturbed quadratic and linear parts are invariant under every
coordinate permutation. If the perturbed objective is preserved on the
feasible set, its coefficient difference is constant within each block;
this necessary condition follows by varying coordinates within their
balance hyperplane. Consequently the coefficient multiset of each block
must equal the coefficient multiset of its image block up to translation.
Put $c=\eta/(3t\,n^2)$ and $\ell=3t$. The sorted successive differences
in block $b$ are
\[
 c\bigl(2((b-1)\ell+i)+1\bigr),\qquad i=1,\ldots,\ell-1.
\]
Their first entries distinguish all blocks. Translation preserves these
differences, so every block is fixed. Within a fixed block, comparison
of the minimum shows that a finite coefficient set cannot equal a
nonzero translate of itself. The translation is zero, and its distinct
entries force the identity permutation.
\end{proof}
Distinctness alone would not prove the statement modulo several balances:
two blocks with coefficient lists differing by additive offsets could be
exchanged without changing the objective on their equal-demand slices.
The squared-index choice excludes this. The symmetry assertion concerns
the original optimization problem; redundant auxiliary variables can
introduce symmetries of their own.

\subsection{A short alternative certificate and the scope of the obstruction}
The blocks have fixed dimension. If component certificates may be
combined by adding their bounds, solving each block separately gives a
certificate of size proportional to $G$; even spatial certification of
each block by Proposition~\ref{prop:chord-upper} has fixed cost for fixed
$t$. Theorem~\ref{thm:relative-cover} counts regions in one global cover
under the stated node oracle and does not constrain such a method. In this
family the explicit sorted optimizer solves each component even more
simply.

A classical coupled cut also closes every root studied here. Its low
written degree does not mean that its nonnegativity follows from a
low-degree local preordering.

\begin{proposition}[A Boolean-quadric clique cut closes the root]
\label{prop:clique-escape}
For integers $1\le k\le n-1$, the inequality
\begin{equation}\label{eq:clique-cut}
 2\sum_{i<j}X_{ij}-2k\sum_i\mu_i+k(k+1)\ge0
\end{equation}
is valid on the continuous quadratic graph
$(\mu,X)=(x,xx^T)$ for $x\in[0,1]^n$.
Adding it to the root equality-product relaxation yields the exact lower
bound $1/4$ for $F$, and the exact root bound for $F_d$ and $C_d$.
One such inequality per block makes the root bound for
\eqref{eq:relative-problem} exact.
\end{proposition}
\begin{proof}
This is Padberg's clique inequality
\citep[Lemma 2, equation (17)]{Padberg1989}, with the whole coordinate
set and parameter $k$, after multiplying by two and changing sides.
Its continuous validity also has a direct proof. After substituting
$X_{ij}=x_ix_j$, the left side is multiaffine. Successively choosing
endpoints can only decrease its value, so its minimum on the cube is
attained at a Boolean vertex. At a vertex with $s$ ones its value is
$(s-k)(s-k-1)\ge0$, since $s-k$ is an integer.

The balance and its linearized products imply
$\sum_{i,j}X_{ij}=K\sum_i\mu_i=K^2$.
Equation~\eqref{eq:clique-cut} therefore gives
\[
 \sum_iX_{ii}\le K^2-2kK+k(k+1)=k+\tfrac14,
 \qquad \sum_i(\mu_i-X_{ii})\ge\tfrac14.
\]
For the perturbed objective, first moments meet the unit bounds and
balance, and their linear cost is at least
$\sum_{i\le k}d_i+d_{k+1}/2$, the linear-program minimum proved in
Corollary~\ref{cor:unique-cardinality}. Adding the two bounds gives
$F_d^*$. The actual optimizer supplies an attaining exact moment matrix,
so the root bound is exact. Adding the conserved $K$ treats $C_d$,
and applying the argument separately to all block balances treats $C$.
The deduction does not use PSD. For the optional boundary case $k=0$,
the same polynomial is $2\sum_{i<j}x_ix_j\ge0$ and the algebra still
applies directly.
\end{proof}

Thus these easy-to-optimize families show limits of particular spatial
certificate systems, including full node preorderings at the specified
orders, while known global cuts avoid the obstruction. In a penalty lift
the coupled affine inequality $\sum_i y_i\ge1/4$ would also close the
unperturbed root; its validity is exactly the global conclusion, so it is
not granted as a local graph inequality. Obstructions that survive
such coupled cuts need different families. Section~\ref{sec:xor-quadratic}
gives one that survives all valid quadratic box inequalities imposed on
moments: signed cubic XOR objectives, which are unrelated to the positive
multilinear gap examples of the earlier sections.
